% =========================================================
% zadani.tex -- Domácí úkol: Taylorovy řady
% Seminář IVT, 3. ročník gymnázia
%
% Kompilace: XeLaTeX
%   latexmk -xelatex -cd zadani/zadani.tex
% =========================================================

\documentclass[12pt, a4paper]{article}

% --- Čeština ---
\usepackage{polyglossia}
\setmainlanguage{czech}
\usepackage{fontspec}

% --- Okraje ---
\usepackage[top=20mm, bottom=20mm, left=22mm, right=22mm]{geometry}

% --- Matematika ---
\usepackage{amsmath}

% --- Zdrojový kód ---
\usepackage{xcolor}
\usepackage{listings}

\definecolor{codebg}{RGB}{245,245,245}
\definecolor{codeblue}{RGB}{0,0,180}
\definecolor{codegreen}{RGB}{0,120,0}
\definecolor{codepurple}{RGB}{120,0,140}

\lstdefinestyle{python}{
    language=Python,
    backgroundcolor=\color{codebg},
    basicstyle=\ttfamily\small,
    keywordstyle=\color{codeblue}\bfseries,
    commentstyle=\color{codegreen}\itshape,
    stringstyle=\color{codepurple},
    breaklines=true,
    showstringspaces=false,
    frame=single,
    tabsize=4,
    numbers=none,
}
\lstset{style=python}

% --- České uvozovky ---
\usepackage{csquotes}
\newcommand{\uv}[1]{\enquote{#1}}

% --- Rámečky pro zadání ---
\usepackage{mdframed}
\newmdenv[
    linewidth=1.5pt,
    linecolor=blue!60!black,
    backgroundcolor=blue!3,
    innertopmargin=10pt, innerbottommargin=10pt,
    innerleftmargin=12pt, innerrightmargin=12pt,
]{zadanibox}

\newmdenv[
    linewidth=1.5pt,
    linecolor=orange!70!black,
    backgroundcolor=orange!5,
    innertopmargin=10pt, innerbottommargin=10pt,
    innerleftmargin=12pt, innerrightmargin=12pt,
]{bonusbox}

% --- Hyperlinks ---
\usepackage{hyperref}
\hypersetup{colorlinks=true, linkcolor=blue!60!black}

% =========================================================
\begin{document}
% =========================================================

\begin{center}
    {\LARGE \textbf{Domácí úkol: Taylorovy řady}}\\[0.5em]
    {\large Seminář IVT --- 3.~ročník gymnázia}\\[0.3em]
    {\normalsize Odevzdat do týdne od hodiny (uploadem do Google Classroom)}
\end{center}

\bigskip
\hrule
\bigskip

\section*{Cíl}

Naučit se implementovat \textbf{Taylorovu řadu pro $\sin$ a $\cos$}
a \textbf{Newtonovu iteraci pro odmocninu} v~čistém Pythonu ---
bez použití \texttt{math.sin}, \texttt{math.cos}, \texttt{math.sqrt}
pro samotný výpočet.
Funkci \texttt{math} použijete jen jako \uv{správnou odpověď}
pro porovnání.

\bigskip

\section*{Povinné zadání}

\begin{zadanibox}
Napište soubor \texttt{ukol.py} s~následujícími funkcemi:

\begin{enumerate}

    \item \textbf{\texttt{faktorial(n)}} ---
          vypočítá $n! = 1 \cdot 2 \cdot 3 \cdots n$
          pomocí smyčky (bez \texttt{math.factorial}).

    \item \textbf{\texttt{sin\_taylor(x, n)}} ---
          vrátí aproximaci $\sin(x)$ pomocí prvních $n$~členů řady:
          \[
              \sin(x) \approx
              \sum_{k=0}^{n-1} \frac{(-1)^k\, x^{2k+1}}{(2k+1)!}
          \]
          Argument $x$ je v~\textbf{radiánech}.

    \item \textbf{\texttt{cos\_taylor(x, n)}} ---
          vrátí aproximaci $\cos(x)$ pomocí prvních $n$~členů řady:
          \[
              \cos(x) \approx
              \sum_{k=0}^{n-1} \frac{(-1)^k\, x^{2k}}{(2k)!}
          \]

    \item \textbf{\texttt{sqrt\_newton(a, n)}} ---
          vrátí aproximaci $\sqrt{a}$ po $n$~iteracích:
          \[
              x_0 = a, \qquad
              x_{k+1} = \frac{x_k + a/x_k}{2}
          \]

\end{enumerate}

\medskip

\textbf{Výstup skriptu}\\
Skript musí vytisknout tabulku porovnávající vaše hodnoty
s~hodnotami z~\texttt{math.sin} a \texttt{math.sqrt}
pro \textbf{alespoň tři různé vstupy} a \textbf{alespoň tři
různé počty členů/iterací} --- viz ukázka níže.

\end{zadanibox}

\bigskip

\subsection*{Ukázka očekávaného výstupu}

\begin{lstlisting}
x = pi/6  (30 stupnu = 0.5236 rad)
n=1: sin_taylor = 0.523599   math.sin = 0.500000   chyba = 2.36e-02
n=3: sin_taylor = 0.500002   math.sin = 0.500000   chyba = 2.12e-06
n=5: sin_taylor = 0.500000   math.sin = 0.500000   chyba = 8.83e-11

sqrt(2), n iteraci:
n=1: sqrt_newton = 2.000000  math.sqrt = 1.414214  chyba = 5.86e-01
n=3: sqrt_newton = 1.414216  math.sqrt = 1.414214  chyba = 2.12e-06
n=5: sqrt_newton = 1.414214  math.sqrt = 1.414214  chyba = 1.59e-12
\end{lstlisting}

\bigskip

\section*{Bonusová část}

\begin{bonusbox}
\textbf{Bonus A} --- Převod stupňů na radiány:
napište funkci \texttt{deg\_to\_rad(stupne)}, která převede
stupně na radiány pomocí $x = \mathrm{stupne} \cdot \pi / 180$,
kde $\pi$ můžete zadat jako konstantu (\texttt{3.14159265358979}).

\medskip

\textbf{Bonus B} --- Graf konvergence:
nainstalujte \texttt{matplotlib} a vykreslete graf absolutní chyby
$|\texttt{sin\_taylor}(x, n) - \sin(x)|$ jako funkce $n$
pro $x = \pi/6$. Osa $y$ má být logaritmická.

Návod pro instalaci: \texttt{pip install matplotlib}

\medskip

\textbf{Bonus C} --- Identita a rychlejší konvergence:\\
Platí $\sin(x) = \cos\!\left(\dfrac{\pi}{2} - x\right)$.
Napište funkci \texttt{sin\_pres\_cos(x, n)}, která vypočítá $\sin(x)$
pomocí vaší funkce \texttt{cos\_taylor} s~touto identitou.

Porovnejte přesnost pro $x = 5\pi/6$ (150°) s~$n = 5$ členy:
\begin{itemize}
    \item \texttt{sin\_taylor(5*pi/6,~5)} --- přímý výpočet Taylorovou řadou
    \item \texttt{sin\_pres\_cos(5*pi/6,~5)} --- výpočet přes identitu
\end{itemize}
Který způsob dá přesnější výsledek?
Zamyslete se: argument $\tfrac{\pi}{2} - \tfrac{5\pi}{6} = -\tfrac{\pi}{3}$
je menší než $\tfrac{5\pi}{6}$ --- proč to vede k lepší konvergenci?
\end{bonusbox}

\bigskip
\hrule
\bigskip

\subsection*{Co je povoleno}

\begin{itemize}
    \item \texttt{import math} --- ale jen pro \texttt{math.sin},
          \texttt{math.cos}, \texttt{math.sqrt}, \texttt{math.pi}
          jako \emph{referenční hodnoty} pro porovnání,
          ne pro samotný výpočet.
    \item Libovolné cykly, podmínky, funkce.
    \item Pro bonus B: \texttt{import matplotlib.pyplot}.
\end{itemize}

\subsection*{Hodnocení}

\begin{itemize}
    \item Povinné zadání: \textbf{8~bodů} (2 body za každou funkci)
    \item Správný formátovaný výstup: \textbf{2~body}
    \item Bonus A (převod stupňů): \textbf{1~bod}
    \item Bonus B (graf): \textbf{2~body}
    \item Bonus C (identita): \textbf{2~body}
\end{itemize}

\end{document}
